\section{Environment and Compatibility Challenges}
\label{sec:env_challenges}

Implementing a fully autonomous MLOps pipeline for embedded systems requires navigating a complex landscape of software dependencies, legacy constraints, and proprietary toolchain limitations. This section details the specific architectural solutions adopted to address these interoperability challenges.

\subsection{Dual-Environment Strategy for Legacy Model Support}
A critical challenge encountered during the development of the framework was the incompatibility between modern Deep Learning libraries required for the agentic orchestration and the legacy formats often found in embedded AI repositories.

\begin{itemize}
    \item \textbf{The Conflict}: The NNI optimization engine and the LangGraph orchestrator rely on modern Python environments (e.g., Keras 3.x, PyTorch 2.x) to leverage recent features in varying generic code generation and asynchronous execution. However, a significant portion of pre-trained embedded AI models (e.g., from older STModelZoo releases or academic papers) are serialized in older formats (Keras 2.x H5) which are essentially incompatible with Keras 3.x execution runtimes due to breaking changes in the serialization API.
    
    \item \textbf{The Solution}: To resolve this, I implemented a \textbf{Dual-Environment Architecture}. The system maintains two distinct virtual environments:
    \begin{enumerate}
        \item \textbf{Primary Environment (\texttt{stm32\_agent})}: Hosting the LLM agents, LangGraph orchestrator, and NNI manager. This environment uses the latest libraries for maximum performance and security.
        \item \textbf{Legacy Environment (\texttt{stm32\_legacy})}: A strictly pinned environment (TensorFlow 2.15, Keras 2.15) dedicated solely to the introspection and analysis of older model files.
    \end{enumerate}
    
    \item \textbf{Mechanism}: When the \textit{Model Discovery Agent} identifies a potential candidate model, it does not attempt to load it directly. Instead, it delegates the inspection task to a subprocess running in the \texttt{stm32\_legacy} environment. This subprocess loads the model, extracts its architecture and I/O signature, and returns the metadata as a JSON object to the main orchestrator, effectively bridging the "semantic gap" between software generations without compromising the stability of the main application.
\end{itemize}

\subsection{Model Format Standardization: H5 vs. TFLite}
While the TensorFlow Lite (\texttt{.tflite}) format is the de facto standard for mobile deployment, the specific requirements of the STEdgeAI Core CLI introduced constraints that necessitated a standardized approach to model persistence.

Experimental validation revealed that the STEdgeAI CLI exhibits higher reliability and broader quantization support when processing Keras H5 files compared to TFLite flatbuffers, particularly for complex custom architectures. Specifically, the CLI's internal graph optimization passes were found to occasionally drop quantization parameters when ingesting pre-quantized TFLite models, leading to performance degradation on the target MCU.

Consequently, the pipeline was engineered to enforce the \textbf{H5 format as the intermediate representation standard}:
\begin{enumerate}
    \item \textbf{NNI Optimization}: The \texttt{trial.py} templates generated by the agents are constrained to save the best-performing models as \texttt{.h5} files, even if the search space includes quantization-aware training.
    \item \textbf{Conversion Pipeline}: The conversion to the final deployable C-code is performed directly from these high-precision \texttt{.h5} artifacts using the STEdgeAI CLI, delegating the quantization layout decisions to the ST toolchain rather than the training framework. This ensures that the generated \texttt{network.c} and \texttt{network.h} files are fully compatible with the X-CUBE-AI runtime engine.
\end{enumerate}
